\documentclass[11pt]{article}
\usepackage[margin=1in]{geometry}
\usepackage{amsmath,amssymb,mathtools,amsthm}
\usepackage{booktabs,array,longtable,multirow}
\usepackage{enumitem}
\usepackage{hyperref}
\usepackage{xcolor}

\newcommand{\SB}{\textsc{ScaleBridge}}
\newcommand{\R}{\mathbb{R}}
\newcommand{\Bun}{\mathcal{B}}
\newcommand{\Run}{\mathcal{R}}
\newcommand{\Lineage}{\mathcal{L}}
\newcommand{\Exec}{\mathcal{E}}
\newcommand{\Norm}{\mathcal{N}}
\newcommand{\InputSpec}{\mathcal{U}}
\newcommand{\StateSpec}{\mathcal{S}}
\newcommand{\Disc}{\mathcal{D}}
\newcommand{\PH}{\mathcal{P}}
\newcommand{\Version}{\mathcal{V}}
\newcommand{\Diag}{\Pi}
\newtheorem{invariant}{Invariant}
\newtheorem{proposition}{Proposition}

\title{\SB{} Phase E0 --- E0-7\\
Runtime, PHVAC Reconstruction, and Portable Model Bundles\\
Mathematical and Scientific Contract v1.1}
\author{ScaleBridge Research Framework}
\date{Ratified August 2026}

\begin{document}
\maketitle

\begin{abstract}
This document freezes the E0-7 mathematical and scientific contract for turning
an already trained or estimated thermal model into an immutable, portable,
forward-executable model artifact. E0-7 is deliberately post-training and
post-estimation: it does not own gradient descent, inverse-PINN training,
CasADi/IPOPT parameter estimation, Bayesian inference, hyperparameter tuning,
MPC, reinforcement learning, or the future SmartBuildingsSim environment.
Instead it defines the common deployment envelope that all later estimation
families must export into and the common forward-runtime contract that Sim1,
Sim2, Sim3, future simulators, MPC, and RL can consume.

The E0-7 bundle contains all static information required for a scientifically
correct forward execution: model identity, Phase-B/Phase-C/Phase-D lineage and
portable data locators, the method-specific executable fitted payload, control
and disturbance input schemas, fitted normalization and inverse-normalization
metadata where required, state/initialization semantics, discretization/runtime
metadata, PHVAC artifacts and reconstruction metadata, and minimal
schema/integrity information. Time-varying forcing trajectories and mutable
runtime state are not part of the immutable model artifact.

For the current ScaleBridge Phase-C contract, each applicable aggregate-zone
PHVAC model is trained as
\[
  \hat p_i=g_{\phi_i}(|Q_{AC,i}|),
\]
against the equal allocated target
\[
  y_i=\frac{P_{\mathrm{HVAC,facility}}}{N},
\]
where $N$ is the total aggregate-zone count. Let $M$ denote the number of
aggregate zones without an applicable PHVAC model; therefore $N-M$ PHVAC models
are available. E0-7 reconstructs the facility-level prediction by
\[
\boxed{
\hat P_{\mathrm{HVAC,facility}}=
\begin{cases}
\displaystyle\sum_{i=1}^{N}\hat p_i, & M=0,\\[6pt]
\displaystyle\frac{N}{N-M}\sum_{i\in\mathcal H}\hat p_i, & 0<M<N,\\[8pt]
\text{unavailable}, & M=N,
\end{cases}}
\]
where $\mathcal H$ is the set of zones with applicable PHVAC models and
$|\mathcal H|=N-M$. The completion term is an allocation correction induced by
the Phase-C training target; it must not be interpreted as HVAC electricity
consumed by a structurally non-HVAC zone. For $M=0$, all $N$ PHVAC models are
available and ordinary summation is used with no completion. For $M=N$, no
model-based PHVAC reconstruction is available.
\end{abstract}

% =============================================================================
\section{Scope and responsibility boundary}

Thermal RC state-space modelling and numerical propagation are already owned by
E0-3 through E0-6. Thermal-network construction is grounded in standard RC
building models \cite{BacherMadsen2011HeatDynamics,GhiausAhmad2020ThermalCircuits},
but the packaging, lineage, runtime, replay, and PHVAC-completion rules in this
document are \SB{}-specific design contracts.

E0-7 begins only after a method has produced a final fitted model. The four later
estimation-method phases may obtain that model through classical machine
learning, scientific machine learning, deterministic nonlinear optimization, or
Bayesian/probabilistic estimation. E0-7 receives the final deployable result and
turns it into a portable artifact with one common forward interface.

\begin{invariant}[Post-estimation boundary]
E0-7 does not own parameter estimation, backward differentiation, optimizer
steps, posterior sampling, EM iterations, hyperparameter optimization, model
selection, MPC, RL, or simulator dynamics. E0-7 owns packaging and forward
execution of an already fitted model.
\end{invariant}

\paragraph{Downstream use.}
The same E0-7 bundle may later be used by historical Sim1/Sim2/Sim3 evaluators,
a future SmartBuildingsSim Gymnasium environment, or other forward consumers.
Those consumers own the source of time-varying controls and disturbances.

% =============================================================================
\section{Immutable portable model artifact}

The canonical immutable artifact is
\[
\boxed{
\Bun=
(\mathcal I,\Lineage,\Exec,\InputSpec,\Norm,\StateSpec,\Disc,\PH,\Version).
}
\]
The components are:
\begin{align*}
\mathcal I &: \text{scientific and model identity},\\
\Lineage &: \text{portable Phase-B/C/D lineage and data locators},\\
\Exec &: \text{method-specific fitted executable payload},\\
\InputSpec &: \text{control/disturbance input schema},\\
\Norm &: \text{fitted normalization and inverse-normalization contract},\\
\StateSpec &: \text{state and initialization semantics},\\
\Disc &: \text{forward solver/discretization contract},\\
\PH &: \text{PHVAC models, transforms, applicability, and reconstruction metadata},\\
\Version &: \text{minimal schema/version/integrity metadata}.
\end{align*}

\begin{invariant}[Forward sufficiency]
The immutable artifact contains every static quantity required to reconstruct and
execute the fitted model correctly in forward mode. It does not contain a
particular simulation's mutable state or its time-varying forcing trajectory.
\end{invariant}

% =============================================================================
\section{Immutable model versus mutable runtime instance}

A loaded model bundle is not itself a running simulation. Let
\[
\boxed{\Run_k=(\Bun,S_k)}
\]
denote a runtime instance, where $S_k$ is mutable state at step $k$.
Examples include
\[
S_k=X_k \quad \text{for an RC model},
\]
\[
S_k=z_k \quad \text{for a Neural ODE latent state},
\]
and
\[
S_k=(h_k,c_k,\text{required history})
\quad \text{for recurrent neural models}.
\]
For a probabilistic runtime, $S_k$ may additionally contain a belief state if
the deployed method explicitly requires one.

\begin{invariant}[No mutable-state serialization as model authority]
The portable scientific model is immutable. A particular running state is
created by initialization and subsequent controls/disturbances and is therefore
owned by the runtime instance, not by the bundle authority.
\end{invariant}

% =============================================================================
\section{Portable Phase-B/C/D lineage and data resolution}

The bundle stores logical locators rather than machine-specific absolute paths.
A generic locator is
\[
\ell=(s,c,r,i,a,w,z,k,p,h),
\]
where $s$ identifies source stage B, C, or D; $c$ a campaign; $r$ an
authoritative run; $i$ a case; $a$ aggregation identity; $w$ weighting policy;
$z$ aggregate-zone identity; $k$ artifact kind; $p$ a path relative to a
registered data root; and $h$ an optional content/schema hash.

The lineage set is
\[
\boxed{\Lineage=(\ell_B,\ell_C,\ell_D)}
\]
with entries present as scientifically applicable. A machine-specific resolver
\[
\mathfrak R_m:\ell\mapsto \text{local filesystem/resource path}
\]
combines the registered local root on machine $m$ with the stored relative
locator.

\begin{invariant}[Portable lineage]
Absolute laptop/home-PC/lab-PC/Kamiak paths are not scientific authority.
Campaign/run/case/aggregation identities plus relative locators and hashes are
the portable authority.
\end{invariant}

\paragraph{Why retain B, C, and D.}
Phase B preserves aggregation provenance; Phase C preserves QAC/PHVAC model
provenance; Phase D identifies the canonical thermal-model dataset used for the
particular training/estimation run. The bundle references these artifacts but
does not duplicate entire historical datasets.

% =============================================================================
\section{Control and disturbance input contract}

Every fitted model exposes a generic forward interface in terms of model state,
controllable inputs, and exogenous disturbances:
\[
\boxed{
(S_{k+1},Y_{k+1})=F_{\Bun}(S_k,C_k,D_k).
}
\]
The exact members of $C_k$ and $D_k$ are declared by $\InputSpec$. For a current
thermal RC model, $C_k$ may include $Q_{AC,k}$, while $D_k$ can contain a model-
applicable subset of
\[
T_{o,k},\;Q_{ZIC,k},\;Q_{ZIR,k},\;Q_{Sol1,k},\;Q_{Sol2,k},\ldots
\]
with E0-2/E0-4 semantics retained.

The model does not own the generation of $C_k$ and $D_k$. A source resolver may
produce
\[
(C_k,D_k)=
\begin{cases}
(C_k^{\mathrm{hist}},D_k^{\mathrm{hist}}), & \text{historical replay},\\
(C_k^{\mathrm{sim}},D_k^{\mathrm{sim}}), & \text{future simulator},\\
(C_k^{\mathrm{caller}},D_k^{\mathrm{caller}}), & \text{other explicit caller}.
\end{cases}
\]

\subsection{Historical Sim1/Sim2/Sim3 source semantics}
For historical evaluation, a replay provider resolves the authoritative
Phase-D/C lineage and supplies observations, disturbances, setpoints/context,
and recorded QAC where the selected evaluation protocol requires it. Sim1 and
Sim2 may replay recorded QAC. Sim3 may replay historical exogenous disturbances
while generating its control/QAC through its evaluator/controller contract.

\subsection{Future SmartBuildingsSim source semantics}
The future SmartBuildingsSim layer will generate or expose simulator
controls/disturbances and pass them into the same E0-7 runtime interface. MPC and
RL will operate through that simulator layer, not by reading Phase-D files
directly.

% =============================================================================
\section{Normalization and inverse-normalization}

If a method is trained in normalized coordinates, the fitted normalization is a
part of the immutable model itself. Let
\[
\tilde u=\mathcal N_u(u),\qquad
\tilde y=F_{\Exec}(\tilde u),\qquad
y=\mathcal N_y^{-1}(\tilde y).
\]
The bundle therefore stores every fitted statistic/transform needed to evaluate
$\mathcal N_u$ and $\mathcal N_y^{-1}$ exactly as used during training.

\begin{invariant}[Physical-unit runtime API]
External E0-7 callers use canonical physical units. Any normalization required
by the fitted method is performed internally by the loaded immutable artifact,
and outputs are returned in canonical physical units unless an explicitly
requested diagnostic representation says otherwise.
\end{invariant}

The training/estimation side is responsible for exporting the correct fitted
normalization object/metadata. E0-7 is responsible for ingesting, validating,
and applying it.

% =============================================================================
\section{Generic method-family payload contract}

E0-7 freezes one common outer bundle but intentionally leaves exact payload
schemas extensible until the four later method classes are implemented. Every
$\Exec$ must declare at minimum: method family and method identity, required
forward-state representation, required input schema, required normalization,
required fitted deployment parameters/weights, output schema, and reconstruction
instructions.

\begin{longtable}{p{0.19\textwidth}p{0.31\textwidth}p{0.41\textwidth}}
\toprule
Family / method & Fitted deployment authority & E0-7 forward implication\\
\midrule
\endfirsthead
\toprule
Family / method & Fitted deployment authority & E0-7 forward implication\\
\midrule
\endhead
Classical ML & network architecture, fitted weights, lag/history contract, fitted normalization & Static or recurrent forward prediction in the method-declared state representation.\\
ODE/RC estimation & final physical $\Theta^\star$ plus RC specification & Reconstruct RC forward equations and E0-5 solver; training machinery is not needed.\\
Inverse PINN-RC & final physical $\Theta^\star$ plus RC specification & Deploy the identified RC physics. Training PINN network weights are not part of the forward authority unless a future explicitly different method declares otherwise.\\
Neural ODE (NODE) & learned vector-field architecture and weights $\psi^\star$ & Integrate the learned vector field from its declared state.\\
Base PINODE & physical $\Theta^\star$ and required learned weights $\psi^\star$ & Reconstruct the hybrid physical/neural RHS and numerical propagation.\\
EBP-PINODE & physical $\Theta^\star$, learned weights $\psi^\star$, and required EBP/projector structure & Reconstruct the complete executable projected/hybrid dynamics.\\
Optimization / CasADi-IPOPT & final physical $\Theta^\star$ plus RC/runtime specification & Treat $\Theta^\star$ as the fitted physical model; the IPOPT NLP is estimator provenance, not the deployed model.\\
Bayesian / probabilistic & $\Theta^\star$, or posterior summary such as $(\mu_\Theta,\Sigma_\Theta)$, or a later-defined posterior payload & Use the same RC forward physics under an explicitly declared deterministic or probabilistic deployment policy.\\
\bottomrule
\end{longtable}

\paragraph{Bayesian deployment policy.}
E0-7 does not force all Bayesian methods into one posterior representation. A
method may export a point parameter set $\Theta^\star$, a mean/covariance pair,
or another future posterior object. The method payload must declare how that
object is turned into a forward prediction, for example posterior mean, MAP,
sampled parameter propagation, or posterior predictive propagation.

% =============================================================================
\section{PHVAC auxiliary-model contract}

PHVAC is electrical HVAC power and remains separate from the thermal model's
signed HVAC heat input. Thermal propagation uses $Q_{AC}$; PHVAC evaluation
uses the fitted Phase-C electrical-power model. PHVAC never enters the thermal
RHS.

\subsection{Current Phase-C equal-allocation training contract}
For all current ScaleBridge campaigns, the PHVAC feature transform is
\[
\boxed{x_i=|Q_{AC,i}|}
\]
and each aggregate-zone PHVAC regression is trained against
\[
\boxed{y_i=\frac{P_{\mathrm{HVAC,facility}}}{N}},
\]
where $N$ is the total number of aggregate thermal zones for the case. E0-7
stores/executes the Phase-C feature-transform metadata rather than treating
absolute value as an eternal hard-coded rule, but the current implementation
must reproduce this equal-allocation contract exactly.

Let
\[
\mathcal Z=\{1,\ldots,N\}
\]
be the aggregate thermal-zone set and let
\[
\mathcal H\subseteq\mathcal Z
\]
be the subset with applicable PHVAC models. Define
\[
M=N-|\mathcal H|,
\qquad |\mathcal H|=N-M,
\]
so that $M$ is the number of aggregate zones without an applicable PHVAC model.
For $i\in\mathcal H$,
\[
\hat p_i=g_{\phi_i}(|Q_{AC,i}|).
\]
A structurally non-HVAC zone is not assigned a fabricated PHVAC model.

\subsection{Equal-allocation completion}
If $M=0$, all $N$ PHVAC models are available and the normal reconstruction is
\[
\boxed{
\hat P_{\mathrm{HVAC,facility}}
=\sum_{i=1}^{N}\hat p_i.
}
\]
If $0<M<N$, the raw available sum
\[
S_{N-M}=\sum_{i\in\mathcal H}\hat p_i
\]
represents only an $(N-M)/N$ share of the Phase-C equal-allocation target. Hence
\[
\boxed{
\hat P_{\mathrm{HVAC,facility}}
=\frac{N}{N-M}S_{N-M},
\qquad 0<M<N.
}
\]
Equivalently, with
\[
\bar p_{\mathcal H}=\frac{1}{N-M}S_{N-M},
\]
the allocation-completion term is
\[
\boxed{
P_{\mathrm{completion}}
=M\bar p_{\mathcal H}
}
\]
and
\[
\hat P_{\mathrm{HVAC,facility}}
=S_{N-M}+P_{\mathrm{completion}}.
\]

\begin{invariant}[Interpretation of completion]
$P_{\mathrm{completion}}$ is an arithmetic reconstruction of the unrepresented
Phase-C equal-allocation target. It is not attributed to a structurally
non-HVAC zone and must not be labelled as that zone's electrical consumption.
\end{invariant}

\subsection{The boundary contracts $M=0$ and $M=N$}
For $M=0$, all $N$ PHVAC models are available; E0-7 uses ordinary summation and
performs no numerical completion. For $M=N$, no PHVAC model is available and
\[
\boxed{M=N\Longrightarrow
\hat P_{\mathrm{HVAC,facility}}^{\mathrm{model}}=\text{unavailable}.}
\]
E0-7 never evaluates $N/(N-M)$ at $M=N$ and never fabricates model-based
facility PHVAC. Historical measured facility HVAC electricity may still be
recorded separately as ground truth when available.

\subsection{Future generalized allocation identity}
For mathematical completeness only, suppose a future Phase-C contract trained
zone $i$ against
\[
p_i=a_iP_{\mathrm{facility}},
\qquad a_i\ge0,
\qquad \sum_{i=1}^{N}a_i=1.
\]
Then for available set $\mathcal H$ with
\[
A_{\mathcal H}=\sum_{i\in\mathcal H}a_i>0,
\]
a natural allocation reconstruction is
\[
\boxed{
\hat P_{\mathrm{facility}}
=\frac{\sum_{i\in\mathcal H}\hat p_i}{A_{\mathcal H}}.
}
\]
The current production implementation does not require arbitrary $a_i$ because
all current ScaleBridge PHVAC campaigns use equal allocation $a_i=1/N$.

% =============================================================================
\section{Forward output and optional diagnostics}

The minimal high-speed step result is
\[
\boxed{
O_k^{\mathrm{basic}}=(S_{k+1},Y_{k+1},P_{\mathrm{HVAC},k}).
}
\]
When validation/evaluation diagnostics are requested, E0-7 may additionally
return
\[
\boxed{
O_k^{\mathrm{diag}}=
(S_{k+1},Y_{k+1},C_k^{\mathrm{used}},D_k^{\mathrm{used}},
\{\hat p_{i,k}\}_{i\in\mathcal H},P_{\mathrm{completion},k},
P_{\mathrm{HVAC},k},\Diag_k).
}
\]
The diagnostic provenance $\Diag_k$ can record the resolved historical/simulator
source, artifact IDs, PHVAC applicability, and reconstruction policy.

\begin{invariant}[Optional validation cost]
Validation capability is mandatory, but expensive integrity/replay checks and
expanded disturbance/provenance echo are optional during ordinary forward
runtime.
\end{invariant}

% =============================================================================
\section{Sim1, Sim2, and Sim3 use of E0-7}

E0-7 provides primitives; the evaluator defines the evaluation protocol.

\paragraph{Sim1.}
Observed/test state is used to initialize/reset the model at each one-step
prediction window and recorded historical disturbances and QAC are resolved from
the authoritative replay source. The model performs only the forward step.

\paragraph{Sim2.}
The model is initialized once, then recursively evolves using its own predicted
state with recorded historical QAC and disturbances.

\paragraph{Sim3.}
The model is initialized once, recursively evolves, replays historical exogenous
disturbances, and receives QAC generated by the Sim3 controller contract.

For every Sim1/Sim2/Sim3 run, the evaluation/recorder must store model-derived
PHVAC results wherever model-based PHVAC is available. For $0<M<N$, the stored
facility result includes the allocation completion; for $M=0$, all $N$ model
predictions are summed directly; for $M=N$, model-based PHVAC is marked
unavailable rather than fabricated. Historical measured PHVAC may be
stored separately as ground truth.

% =============================================================================
\section{Serialization and save/load scientific equivalence}

Let
\[
\operatorname{Save}:\Bun\rightarrow b
\]
serialize a bundle and
\[
\operatorname{Load}:b\rightarrow\Bun'
\]
load it. E0-7 requires scientific equivalence
\[
\boxed{\Bun'\equiv_{\mathrm{sci}}\Bun.}
\]
For physics-based models,
\[
\Theta_{\mathrm{loaded}}^\star
=\Theta_{\mathrm{saved}}^\star
\]
to serialization precision. For required neural deployment weights,
\[
\psi_{\mathrm{loaded}}^\star
=\psi_{\mathrm{saved}}^\star
\]
to serialization precision. Fitted normalization parameters must likewise be
preserved.

Given identical initialization and input sequence,
\[
(S_0,C_{0:K},D_{0:K}),
\]
pre-save and post-load trajectories must satisfy
\[
Y_{0:K}^{\mathrm{pre}}
\approx
Y_{0:K}^{\mathrm{post}},
\]
and, whenever available,
\[
P_{\mathrm{HVAC},0:K}^{\mathrm{pre}}
\approx
P_{\mathrm{HVAC},0:K}^{\mathrm{post}}
\]
within declared dtype/backend tolerances.

\paragraph{Portable specification versus backend cache.}
The authoritative portable representation is a versioned scientific
specification plus standard numerical/model payloads. A backend-specific compiled
object may be cached for performance but is reconstructable convenience, not the
scientific authority. E0-7 must not make a machine-specific Python pickle the
only representation of a scientific model.

% =============================================================================
\section{Deterministic replay}

For a deterministic bundle, fixed initialization, fixed input sequence, and
fixed deterministic runtime configuration, repeated execution must reproduce the
same trajectory within declared numerical tolerances:
\[
F_{\Bun}(S_0,C_{0:K},D_{0:K})^{(1)}
\approx
F_{\Bun}(S_0,C_{0:K},D_{0:K})^{(2)}.
\]
This includes save/load replay equivalence. Comprehensive execution of these
invariants across representative methods, flavours, data organizations,
backends, and machines belongs to the E0-10 testing campaign.

% =============================================================================
\section{Minimal schema/version/integrity contract}

To avoid unnecessary infrastructure, E0-7 requires only:
\[
v_{\mathrm{bundle}},
\qquad
v_{\mathrm{payload}},
\qquad
H_{\mathrm{manifest}}.
\]
A loader accepts supported versions and fails explicitly on an incompatible
unknown version. E0-7 does not require a general migration framework at this
stage.

% =============================================================================
\section{Worked numerical example}

Consider $N=5$ aggregate thermal zones
\[
\mathcal Z=\{A,B,C,D,E\}
\]
and suppose zone $E$ is structurally non-HVAC. Hence
\[
\mathcal H=\{A,B,C,D\},
\qquad M=1,
\qquad |\mathcal H|=N-M=4.
\]
Assume a fitted physics-based model bundle has already stored final physical
parameters $\Theta^\star$, its state ordering, E0-5 runtime solver metadata,
Phase-D historical data locator, fitted normalization if applicable, and four
Phase-C PHVAC models.

At step $k$, historical replay resolves
\[
D_k=(T_{o,k},Q_{ZIC,k},Q_{ZIR,k},Q_{Sol1,k},Q_{Sol2,k},\ldots)
\]
and the evaluation protocol supplies or generates $Q_{AC,k}$ as appropriate.
The thermal model computes
\[
(S_{k+1},Y_{k+1})=F_{\Bun}(S_k,C_k,D_k).
\]

Suppose the four available Phase-C PHVAC predictions are
\[
\hat p_A=950\ \mathrm{W},\quad
\hat p_B=1020\ \mathrm{W},\quad
\hat p_C=980\ \mathrm{W},\quad
\hat p_D=1050\ \mathrm{W}.
\]
Then
\[
S_{N-M}=950+1020+980+1050=4000\ \mathrm{W}
\]
and
\[
\bar p_{\mathcal H}=\frac{4000}{N-M}=\frac{4000}{4}=1000\ \mathrm{W}.
\]
Because current Phase-C training allocates facility PHVAC equally across all
$N=5$ aggregate zones, the unrepresented allocated share is
\[
P_{\mathrm{completion}}
=M\bar p_{\mathcal H}=(1)(1000)=1000\ \mathrm{W}.
\]
Therefore
\[
\boxed{
\hat P_{\mathrm{HVAC,facility}}
=4000+1000
=5000\ \mathrm{W}.
}
\]
Equivalently,
\[
\frac{N}{N-M}S_{N-M}=\frac{5}{4}(4000)=5000\ \mathrm{W}.
\]
The added $1000$ W is not assigned to zone $E$; it is explicitly labelled
allocation completion.

If instead $M=0$, then all five PHVAC models are available and E0-7 simply
sums the five predictions with no allocation completion. At the opposite
boundary, if $M=N=5$, no PHVAC model is available and E0-7 returns
\[
P_{\mathrm{HVAC,facility}}^{\mathrm{model}}=\text{unavailable}
\]
without evaluating $N/(N-M)$. Historical measured facility PHVAC, if present,
remains a separate ground-truth series.

For optional diagnostic output, the step can also return the resolved
$D_k^{\mathrm{used}}$, QAC/control values actually consumed, the four local
PHVAC predictions, the $1000$ W completion term, and the Phase-B/C/D artifact
IDs used for historical replay.

Finally, after serializing and reloading the model bundle, the same
initialization and replay sequence must reproduce the same fitted parameters,
normalization, forward thermal trajectory, and PHVAC reconstruction within the
frozen numerical tolerances.

% =============================================================================
\section{Ratified E0-7 decisions}

The following are frozen by this contract.
\begin{enumerate}[leftmargin=1.8em]
\item The immutable artifact contains all static information needed for correct forward execution.
\item Fitted normalization/inverse-normalization required by a method is part of the immutable artifact.
\item Runtime state is separate from the immutable model.
\item The model exposes generic control and disturbance inputs; historical replay and future simulators own the source of those values.
\item Phase-B/C/D lineage is stored through portable logical locators, not machine-specific absolute paths.
\item Time-varying forcing trajectories are referenced/resolved and are not embedded into every model artifact.
\item Disturbances actually used can be echoed through optional diagnostic/evaluation output.
\item E0-7 keeps exact method payload schemas generic until the later four estimation classes are implemented.
\item One common outer bundle/runtime envelope spans Classical ML, SciML, deterministic optimization, and Bayesian/probabilistic models.
\item ODE/RC and Inverse-PINN deployment use final physical $\Theta^\star$ rather than training-network weights; NODE uses its learned vector field; Base-PINODE and EBP-PINODE retain the fitted physical/neural executable components they require.
\item CasADi/IPOPT output is the fitted physical $\Theta^\star$ plus RC/runtime specification; the estimation NLP is not the deployed scientific model.
\item Bayesian payloads may contain point $\Theta^\star$, mean/covariance, or a later-defined posterior object under an explicit deployment policy.
\item Runtime APIs remain in physical units; normalization is internal to the artifact.
\item Current Phase-C PHVAC uses $|Q_{AC}|$ as its predictor transform and facility-HVAC-electricity divided equally by total aggregate-zone count $N$ as its target allocation.
\item For current equal allocation and $0<M<N$, facility PHVAC is reconstructed by $[N/(N-M)]\sum_{i\in\mathcal H}\hat p_i$; the added quantity is labelled allocation completion.
\item For $M=0$, ordinary summation across all $N$ PHVAC models is used; for $M=N$, model-based PHVAC is unavailable and $N/(N-M)$ is never evaluated.
\item Structural PHVAC absence, valid zero prediction, missing-required artifact, and allocation completion are distinct states.
\item Save/load equivalence, hashes, expanded disturbance echo, normalization round-trip, and detailed replay checks are mandatory validation capabilities but optional during ordinary runtime.
\item E0-7 supplies reusable forward runtime primitives for Sim1/Sim2/Sim3 and later SmartBuildingsSim but does not implement evaluator, simulator, MPC, or RL logic.
\item Every Sim1/Sim2/Sim3 evaluation records model-derived PHVAC wherever available, including building-level reconstruction metadata; historical measured PHVAC may be retained separately as truth.
\item E0-10 is the authoritative comprehensive real E0 testing campaign that exercises these contracts after the four estimation-method families are implemented.
\end{enumerate}

% =============================================================================
\section{Qualification boundary before implementation}

The E0-7 implementation may begin only after this mathematical authority is
installed and validated. Implementation qualification must at minimum test:
\begin{enumerate}[leftmargin=1.8em]
\item bundle schema construction and strict required-field validation;
\item portable B/C/D locator round-trip independent of one machine's absolute path;
\item immutable bundle versus mutable runtime-state separation;
\item physical-unit input/output with normalization round-trip where present;
\item method-payload envelope extensibility without prematurely encoding all E.1--E.4 internals;
\item historical replay input resolution and explicit caller/simulator input compatibility;
\item PHVAC $M=0$, $0<M<N$, and $M=N$ branches;
\item explicit allocation-completion provenance and non-attribution to a non-HVAC zone;
\item save/load parameter/weight/normalization preservation;
\item pre-save versus post-load forward replay equivalence;
\item optional versus normal-runtime diagnostics behavior;
\item no regression of frozen E0-2 through E0-6 contracts.
\end{enumerate}

\bibliographystyle{plain}
\bibliography{ScaleBridge_Thermal_Modeling_References}

\end{document}
